% TIME 2026 — Extended Abstract submission
%
% Target venue: 33rd International Symposium on Temporal Representation
% and Reasoning (TIME 2026), University College Cork, 1–3 September 2026.
% Category: Extended abstract (4 pages, excluding references and appendix).
% Format:   OASIcs (oasics-v2021).
%
% Build:
%   1. Download oasics-v2021 from
%      https://submission.dagstuhl.de/documentation/authors
%      and place oasics-v2021.cls (plus its supporting files) in this
%      directory.
%   2. pdflatex extended-abstract && pdflatex extended-abstract
%
% This file uses an inline thebibliography environment so no .bib is
% required at draft stage. Switch to BibTeX (plainurl) before final
% submission.

\documentclass[a4paper,USenglish,cleveref,autoref,thm-restate]{oasics-v2021}

\title{Time as Intervals, Not Instants: \\
  A Practical Implementation of Allen's Interval Algebra Over ISO~8601}
\titlerunning{Time as Intervals, Not Instants}

\author{Kip Cole}{Independent}{kipcole9@gmail.com}{}{}

\authorrunning{K.~Cole}
\Copyright{Kip Cole}

\ccsdesc[500]{Information systems~Temporal data}
\ccsdesc[300]{Software and its engineering~Domain specific languages}
\ccsdesc[300]{Theory of computation~Modal and temporal logics}

\keywords{Temporal intervals, Allen's interval algebra, ISO 8601,
  calendrical systems, partial dates, set operations on time,
  domain-specific languages}

\relatedversion{Source repository: \url{https://github.com/kipcole9/tempo}}

\begin{document}

\maketitle

\begin{abstract}
  Mainstream programming languages model temporal values as scalar
  instants: a \texttt{Date} is a point, a \texttt{DateTime} is a point,
  and a partial value such as ``June 2026'' has no canonical
  representation at all. Yet the temporal values that humans, calendars,
  and business logic actually manipulate are spans on the timeline ---
  the year 2026, the afternoon of the 15th, an archaeological period,
  a meeting. Tempo is an Elixir library that reverses the polarity:
  every temporal value is a bounded half-open interval at a declared
  resolution, set-theoretic combination is closed under that type, and
  comparison is defined by Allen's thirteen relations rather than by
  scalar ordering. The result eliminates an entire class of off-by-one
  and ``end-of-day'' bugs by construction, gives partial ISO~8601
  values (\texttt{2026}, \texttt{2026-06}, \texttt{2026-06-XX}) an
  unambiguous denotational semantics, and lifts Allen's algebra from a
  textbook abstraction to an everyday predicate vocabulary
  (\texttt{overlaps?}, \texttt{within?}, \texttt{adjacent?}). This
  extended abstract describes the conceptual model, the unification of
  partial dates with explicit intervals, and the design of set
  operations over intervals and interval sets that is the library's
  current focus.
\end{abstract}

\section{Motivation: the instant--span impedance mismatch}

Every mainstream temporal library --- Python's \texttt{datetime}, Java's
\texttt{java.time}, JavaScript's \texttt{Temporal} proposal, Elixir's
own \texttt{Date}/\texttt{Time}/\texttt{DateTime} --- shares a design
choice that goes largely unexamined: a temporal value is a
\emph{scalar instant} on the timeline. A day is the midnight at its
start, a month is the midnight on its first, a year is the midnight
on January~1. The span that humans actually mean when they say ``June
2026'' is materialised, if at all, by ad-hoc helpers
(\texttt{start\_of\_month}, \texttt{end\_of\_day}) that each project
re-implements with its own conventions about inclusivity, time zones,
and DST.

This instant-first model has three persistent consequences:

\paragraph*{Ambiguity at the type level.} The expression
\texttt{Date.new(2026, 6, 15)} denotes a 24-hour span in business
prose and an instant in the type system. The disagreement is resolved
informally, per call site, and is a recurring source of off-by-one
bugs at month, year, and DST boundaries. Tiwari et al.'s recent
empirical study of 151 date/time bugs in open-source Python projects
identifies time-zone-related mistakes as the largest single category,
with most bugs traceable to misconceptions about library API
conventions and edge-case behaviour~\cite{tiwari2025time}; Liva's
earlier work on Java programs reaches a structurally similar
conclusion~\cite{liva2018}.

\paragraph*{No representation for partial dates.} A user who writes
``2022'' or ``the 1560s'' or ``approximately June 2026'' has produced
a perfectly meaningful temporal value, but the type system has no
home for it. Solutions range from sentinel values
(\texttt{Date.new(2022, 1, 1)} stands in for ``the year'') to parallel
type hierarchies (EDTF~\cite{edtf}) that exist outside the standard
library, or to dedicated calendric type systems such as Bry and
Spranger's CaTTS for Semantic Web query
languages~\cite{bry-spranger-multicalendar}. The fragmentation is
structural, not incidental.

\paragraph*{Comparison is reduced to scalar ordering.} Two dates are
either \texttt{<}, \texttt{=}, or \texttt{>}. The richer relational
structure that Allen identified in 1983~\cite{allen1983} ---
\textit{precedes}, \textit{meets}, \textit{overlaps}, \textit{starts},
\textit{during}, \textit{finishes}, \textit{equals}, and their inverses
--- is invisible at the type level, and reimplemented case-by-case in
business logic that interrogates whether a meeting ``overlaps'' the
workday or a contract ``contains'' a billing period.

Tempo's design is the observation that all three problems dissolve
simultaneously if intervals are made the primitive type. There are no
points: an apparent ``instant'' is the interval of one unit at the
finest declared resolution, in the same spirit as the time ontology
$T_{bounded\_meeting}$ of Gr\"uninger and
Li~\cite{gruninger2017allenontology}, in which intervals are the only
entities in the domain and Allen's relations are derivable from a
single primitive (\textit{meets}).

\paragraph*{Empirical coverage.} Against the 151 Python bugs
catalogued by Tiwari et al., the design choices described in the
remainder of this paper address every in-scope bug class: 43 (28\%)
are \emph{avoided by construction} (the na\"ive / aware datetime
split has no Tempo analogue), 56 (37\%) are \emph{directly addressed}
by explicit machinery (zone propagation through set operations, DST
gap rejection at parse time and IXDTF-offset fold disambiguation per
RFC~9557~\S 4.5, calendar-aware arithmetic, half-open materialisation,
microsecond-precision durations), 28 (19\%) are partial, and 24
(16\%) are out of scope (Python deprecations and third-party-library
logic such as holiday tables). \emph{Zero} bugs are classified as
architectural gaps. The per-bug classification is provided as
supplementary material at the repository.

\section{Intervals as the primitive}

Every Tempo value is a bounded half-open interval $[\mathit{first},
\mathit{last})$ at a declared resolution. Internally an interval is
a pair of endpoints, which is the standard finite model of
$T_{bounded\_meeting}$ characterised by Gr\"uninger and
Li~\cite{gruninger2017allenontology} (their Theorem~15): vertices of
the directed graph $\mathfrak{M}^{bounded\_meeting}$ correspond to
intervals, edges to the \textit{meets} relation, and structures up to
isomorphism faithfully capture the ontology's models.

A note on terminology before proceeding: \emph{resolution} as used
throughout this paper is a Tempo-side property of an input value
(\texttt{2026} carries year-resolution, \texttt{2026-06-15} carries
day-resolution) used for partial-date materialisation, default
iteration, and display. It is \emph{not} the granularity of
Bettini, Wang and Jajodia~\cite{bettini-granularities} (a lattice
of partitions of the time line that participates in Allen
comparisons), and it does \emph{not} appear in
$T_{bounded\_meeting}$, where intervals are individuated solely by
their position in the \textit{meets}-structure. Once an interval is
materialised to its $[\mathit{first}, \mathit{last})$ form,
resolution plays no further role in Allen comparisons or in
set-theoretic operations. The word is also distinct from
\emph{precision} and \emph{accuracy} in the measurement-science
sense: Tempo makes no claim about the precision of a value's
construction, only about the extent it denotes.

A second Tempo-side distinction is between \emph{anchored} and
\emph{non-anchored} values. An anchored value carries a year
component (e.g.\ \texttt{2026-06-15} or \texttt{2026}); materialising
it yields an interval that is an \emph{individual} in
$T_{bounded\_meeting}$ in the standard model-theoretic sense ---
a vertex of $\mathfrak{M}^{bounded\_meeting}$ (Theorem~15) with a
definite position in the \textit{meets}-structure. A non-anchored
value (a time-of-day template such as \texttt{T10:30}) is not an
individual: it is an open expression with the date components as
free variables, denoting the family of individuals that result when
those variables are bound. Tempo's \texttt{anchor/2} and
\texttt{align/3} \texttt{:bound} option supply that binding; the
ontological move is \emph{instantiation}, and the result is the
intervalised individual the ontology talks about. The two
operations Tempo exposes for working with non-anchored values are
\emph{projection} (resolving to a specific universal-timeline
extent given an anchor) and \emph{selection} (filtering a set of
candidate anchors). Neither has a counterpart inside
$T_{bounded\_meeting}$ itself, which is silent on schemas; they
sit, like resolution, in the engineering layer above.

The library distinguishes two forms of interval that interconvert
losslessly.

\paragraph*{Implicit spans.} A single ISO~8601 datetime \emph{is} an
interval whose span is determined by the next-higher resolution that
is not stated:
\begin{center}
\begin{tabular}{lcl}
  \texttt{2026}              &$\equiv$& $[\texttt{2026-01-01},\, \texttt{2027-01-01})$ \\
  \texttt{2026-01}           &$\equiv$& $[\texttt{2026-01-01},\, \texttt{2026-02-01})$ \\
  \texttt{2026-01-15}        &$\equiv$& $[\texttt{2026-01-15},\, \texttt{2026-01-16})$ \\
  \texttt{2026-01-15T10}     &$\equiv$& $[\texttt{2026-01-15T10:00},\, \texttt{2026-01-15T11:00})$ \\
\end{tabular}
\end{center}
The partial-date problem disappears: \texttt{2026} is not an
under-specified instant, it is a definite interval of one year (the
distinction is semantic, not a claim about measurement precision or
accuracy). The ``last day of the month'' problem disappears: the
upper bound of \texttt{2026-01} is
\texttt{2026-02-01}, computed from the calendar arithmetic of the
target calendar (Gregorian, Hebrew, Islamic Civil, \dots), and is
correct by construction in leap years and across month-length
variation.

\paragraph*{Explicit spans.} A pair of datetimes written with a range
operator --- \texttt{2026-01-01..2026-02-01} or, in IXDTF-extended
form~\cite{ixdtf}, \texttt{2026-01-01/2026-02-01} --- carries its own
boundaries. Its iteration unit is the resolution of its boundaries,
not of a partial denotation.

\paragraph*{Half-open by design.} The half-open $[\mathit{first},
\mathit{last})$ convention is not notational shorthand; it is enforced
at every level of the API. Adjacent intervals concatenate cleanly:
$[a, b) \cup [b, c) = [a, c)$ with neither overlap nor gap. Iteration
yields a sequence whose union is exactly the parent interval.
Conversion between implicit and explicit forms is a structural
identity, not an arithmetic approximation. Library code that treats
the upper bound as inclusive is, by the design contract, a defect.
Alternative timestamp representations have been explored at TIME ---
notably Dyreson and Ahsan's log-segmented
timestamps~\cite{dyreson2021logsegmented,dyreson2023logsegmented},
which decompose a period into power-of-two segments so that standard
B$^+$-tree indexes can support sequenced and nonsequenced queries on
an unmodified DBMS. The two designs occupy different points in the
same space: log-segmented timestamps optimise query evaluation in a
relational engine; Tempo's half-open intervals optimise expressiveness
and uniformity at the application layer.

\section{Allen's algebra as the predicate vocabulary}

Tempo exposes Allen's thirteen relations~\cite{allen1983} directly as
the result of its core comparison primitive. The first-order ontology
underlying the algebra is characterised by Gr\"uninger and
Li~\cite{gruninger2017allenontology}, who identify the time ontology
logically synonymous with Allen's interval algebra and prove a
representation theorem for its models; Tempo is, in effect, an
ergonomic software realisation of that ontology over the
calendar-aware datetime values defined by ISO~8601. The thirteen
relations are returned directly from the relation primitive:
\begin{verbatim}
  Tempo.relation(a, b) ::
    :precedes | :meets | :overlaps | :finished_by |
    :contains | :starts | :equals |
    :started_by | :during | :finishes | :overlapped_by |
    :met_by | :preceded_by
\end{verbatim}

This is rarely the right surface for application code, however ---
business questions tend to name unions of Allen relations.
``Does the meeting fall within the workday?'' is the set
$\{\textit{equals}, \textit{starts}, \textit{during}, \textit{finishes}\}$;
``does this period touch that one?'' is the complement of
$\{\textit{precedes}, \textit{preceded\_by}\}$. Hand-rolled inline
matching on relation lists is verbose and obscures the intent, so the
library promotes the common unions to first-class predicates ---
\texttt{within?/2}, \texttt{overlaps?/2}, \texttt{adjacent?/2},
\texttt{contains?/2} --- each defined as a specific subset of Allen's
relations and each readable aloud as the English question it answers.

The predicate vocabulary is not decorative. It is a design discipline:
when an example in the cookbook reaches for a relation list or a
manual endpoint comparison, that is taken as evidence that the
abstraction is missing, and the predicate is added before the example
is written. The same discipline applies to durations
(\texttt{at\_least?}, \texttt{shorter\_than?}) and to interval-set
queries (\texttt{any\_overlaps?}, \texttt{disjoint?}). The library's
public surface is, deliberately, an English-prose description of
Allen's algebra over ISO~8601 values.

\section{Position among temporal formalisms}

The choice to make intervals primitive places Tempo on one side of a
long-standing divide. Vilain and Kautz~\cite{vilainkautz1986} take
\emph{points} as primitive: an interval is an ordered pair of endpoint
points, and an interval relation reduces to a conjunction of
point-relations $\{<,=,>\}$ among the four endpoints. The reward is
tractability --- reasoning in the point algebra is polynomial, whereas
constraint reasoning over the full interval algebra is intractable in
general~\cite{vilainkautz1986,nebel1995}. Allen and
Hayes~\cite{allenhayes1989} take the opposite stance, axiomatising the
algebra from the single primitive \textit{meets} and reconstructing
points only as derived entities --- \emph{moments} (intervals with no
proper subinterval) and the meeting-places between intervals.
Gr\"uninger and Li's $T_{bounded\_meeting}$~\cite{gruninger2017allenontology}
continues this interval-only lineage, excluding degenerate intervals
and points from the domain entirely; this is the ontology Tempo
adopts.

Tempo's contribution is best read as a synthesis across this divide.
\emph{Ontologically} it is interval-only: there are no instants,
degenerate intervals are excluded, and the half-open seam realises
Hayes' open-connected model rather than Allen's closed intervals.
\emph{Computationally}, however, comparison projects every endpoint to
a real number on a single UTC frame and orders intervals by their
endpoints --- precisely the Vilain--Kautz reduction, used as an engine
beneath an interval ontology. Two consequences follow. First, because
Tempo's intervals are always \emph{fully grounded} (concrete
endpoints, never disjunctive qualitative constraints), it never enters
the NP-hard regime that motivated the point algebra: relation queries
are constant-time endpoint comparisons rather than constraint
propagation. Second, the ``instant'' is recovered not as a primitive
point but as the interval of one unit at the value's \emph{finest
declared resolution} --- a \emph{moment} in Allen and Hayes' sense,
but relativised to representational precision: a second-resolution
datetime is a one-second interval, a microsecond value a
one-microsecond interval. This resolution-indexed atomicity, tying
interval granularity to ISO~8601 syntax rather than to an absolute
atom, is the element of the model without a direct precedent in the
formalisms it draws on.

\section{Set operations on intervals and interval sets}

The current implementation focus is closing the library under
set-theoretic operations. The intuition is straightforward --- any
two intervals can be combined into a (possibly disjoint, possibly
zero-member) set of intervals --- but the engineering requires careful
handling of the implicit/explicit-span distinction, of cross-calendar
operands, and of cross-zone arithmetic. A single connected, non-empty
result is returned as an \texttt{Interval}; every other shape ---
disjoint, or zero-member --- is returned as an \texttt{IntervalSet}.
An interval is never returned with zero extent, in line with
$T_{bounded\_meeting}$'s exclusion of degenerate intervals from the
domain.

\paragraph*{Pairwise operations.} \texttt{union/2}, \texttt{intersection/2},
and \texttt{difference/2} on a pair of operands $a, b$ return either an
interval or an \texttt{IntervalSet}, depending on whether the result
is connected. Operands may be supplied in either implicit or explicit
form; the library normalises both to the explicit half-open
representation, performs the geometric operation, and re-presents the
result in the most compact form (a single interval when possible).

\paragraph*{$n$-ary operations and canonicalisation.} Extending the
pairwise operations to lists of intervals raises the question of
canonical form. Tempo defines the canonical form of an
\texttt{IntervalSet} as the sorted, coalesced sequence in which no two
intervals overlap or meet. Coalescing is the closure of the union
operation under iteration: $\bigcup_i [a_i, b_i)$ expressed as a
minimal disjoint sequence. This is the form returned by every
set-producing operation in the public API. Coalescing also realises
the Sum Axiom of $T_{bounded\_meeting}$ (axiom~15
in~\cite{gruninger2017allenontology}): three chain-meeting intervals
$\textit{meets}(i,j) \land \textit{meets}(j,k) \land \textit{meets}(k,m)$
are summed into the single interval that runs from $i$'s start to
$m$'s end.

\paragraph*{Cross-calendar and cross-zone semantics.} Because each
operand carries its own calendar and zone, set operations are
defined on the underlying UTC reference frame --- each endpoint is
projected to gregorian seconds (the standard Erlang representation)
so that endpoints from different calendars or zones can be ordered
in a single linear sequence. The word ``instant'' is avoided here
because Gr\"uninger and Li's ontology excludes points from its
domain (intervals are the only entities); naming the projection a
``UTC instant'' would suggest a commitment to a competing
point-based ontology. The projection is purely a computational
device --- a real number used to order endpoints in a single linear
frame. The result's calendar and zone are taken from the first
operand, preserving the caller's intent for downstream display.
This first-operand-wins rule is itself a design choice that deserves
community scrutiny --- alternatives include rejecting heterogeneous
operands, or carrying the heterogeneity through into a sum-typed
result --- and is one of the points the talk would invite feedback
on.

\paragraph*{Daylight-saving transitions.} The UTC projection
interacts with daylight-saving transitions in two distinct ways,
both handled at parse time so set operations on the projected values
need no further disambiguation. A wall time that falls in a DST
\emph{gap} (the non-existent hour at spring-forward) is rejected at
construction with a \texttt{ZoneGapError} --- a stronger guarantee
than addressing the bug, because the malformed value cannot be
constructed. A wall time that falls in a DST \emph{fold} (the
doubled hour at fall-back) is disambiguated by an explicit offset in
the IXDTF suffix, per RFC~9557~\S 4.5: the value
\texttt{2024-11-03T01:30:00-04:00[America/New\_York]} unambiguously
picks the pre-fall-back (EDT) period and
\texttt{2024-11-03T01:30:00-05:00[America/New\_York]} picks the
post-fall-back (EST) period. The disambiguator is part of the
string representation and round-trips through serialisation --- a
property that distinguishes this design from meta-attribute
schemes (such as Python's \texttt{fold} parameter) in which the
choice is carried out-of-band and lost on serialisation.

\paragraph*{Beyond the primitives.} Once \texttt{union},
\texttt{intersection}, and \texttt{difference} are closed, the natural
extensions follow: symmetric difference, containment queries, the
inverse of an interval set against a bounding context (``free time''
given a calendar of meetings), and reductions across an
\texttt{IntervalSet} (total duration, longest gap, densest hour).

\section{Discussion and open questions}

Tempo is a working library; this submission's contribution is a
\emph{design} contribution rather than a \emph{result} contribution,
and the talk would aim to make the design choices legible to an
audience whose research is the theory the library implements. Three
questions seem particularly worth the TIME audience's attention:

\begin{enumerate}
  \item \textbf{Granularity and resolution.} Tempo's
    next-higher-undefined-resolution rule for implicit spans is
    operationally simple but does not reference the more general
    treatment of temporal granularity in the database
    literature~\cite{bettini-granularities} nor the calendric type
    constraints in CaTTS~\cite{bry-spranger-multicalendar}. Is the
    simple rule a sufficient model in practice, or does it accumulate
    edge cases as the calendar set grows (lunar, sidereal, fiscal)?

  \item \textbf{Allen's algebra over interval sets.} The thirteen
    relations, and the $T_{bounded\_meeting}$ ontology that
    underlies them, are defined on convex intervals
    only~\cite{gruninger2017allenontology}. Their lifting to interval
    sets admits several reasonable definitions
    (existential, universal, set-of-relations); the library currently
    offers the existential lifting, but the right default is not
    obvious, and a principled extension of the ontology to non-convex
    interval sets would settle the question.

  \item \textbf{Cross-calendar comparison.} Comparing a Hebrew date
    with a Gregorian one resolves cleanly when both project to the
    same proleptic UTC timeline, but the timeline is itself a
    modelling choice (TAI? UT1? a calendar-specific civil time?).
    Tempo's current choice is pragmatic; a principled framework would
    be a worthwhile contribution from the symposium back to the
    library.
\end{enumerate}

The wider goal is to demonstrate, by existence proof, that the
theoretical structure of temporal reasoning translates into an
ordinary application library --- one that an engineer writes
\texttt{Tempo.overlaps?(a, b)} into without thinking of it as a
research artefact. The talk would walk through the design, the
choices made, the choices still open, and invite the symposium to
sharpen them.

\begin{thebibliography}{9}

\bibitem{allen1983}
  James F. Allen.
  \newblock Maintaining knowledge about temporal intervals.
  \newblock \emph{Communications of the ACM}, 26(11):832--843, 1983.
  \newblock \doi{10.1145/182.358434}

\bibitem{vilainkautz1986}
  Marc B. Vilain and Henry A. Kautz.
  \newblock Constraint propagation algorithms for temporal reasoning.
  \newblock In \emph{Proceedings of the 5th National Conference on
  Artificial Intelligence (AAAI-86)}, pages 377--382, 1986.

\bibitem{allenhayes1989}
  James F. Allen and Patrick J. Hayes.
  \newblock Moments and points in an interval-based temporal logic.
  \newblock \emph{Computational Intelligence}, 5(3):225--238, 1989.
  \newblock \doi{10.1111/j.1467-8640.1989.tb00329.x}

\bibitem{nebel1995}
  Bernhard Nebel and Hans-J\"urgen B\"urckert.
  \newblock Reasoning about temporal relations: A maximal tractable
  subclass of Allen's interval algebra.
  \newblock \emph{Journal of the ACM}, 42(1):43--66, 1995.
  \newblock \doi{10.1145/200836.200848}

\bibitem{iso8601-1}
  International Organization for Standardization.
  \newblock \emph{ISO 8601-1:2019 --- Date and time --- Representations
  for information interchange --- Part 1: Basic rules}, 2019.

\bibitem{iso8601-2}
  International Organization for Standardization.
  \newblock \emph{ISO 8601-2:2019 --- Date and time --- Representations
  for information interchange --- Part 2: Extensions}, 2019.

\bibitem{ixdtf}
  U. Sharma and C. Bormann, editors.
  \newblock \emph{Date and Time on the Internet: IXDTF Extensions}.
  \newblock Internet-Draft draft-ietf-sedate-datetime-extended-09,
  IETF, 2024.
  \newblock \url{https://www.ietf.org/archive/id/draft-ietf-sedate-datetime-extended-09.html}

\bibitem{edtf}
  Library of Congress.
  \newblock \emph{Extended Date/Time Format (EDTF) Specification}, 2019.
  \newblock \url{https://www.loc.gov/standards/datetime/}

\bibitem{bettini-granularities}
  Claudio Bettini, Sushil Jajodia, and Sean X. Wang.
  \newblock \emph{Time Granularities in Databases, Data Mining, and
  Temporal Reasoning}.
  \newblock Springer, 2000.

\bibitem{tiwari2025time}
  Shrey Tiwari, Serena Chen, Alexander Joukov, Peter Vandervelde,
  Ao Li, and Rohan Padhye.
  \newblock It's about time: An empirical study of date and time bugs
  in open-source Python software.
  \newblock In \emph{Proceedings of the 22nd International Conference
  on Mining Software Repositories (MSR '25)}. IEEE/ACM, 2025.
  \newblock Distinguished Paper Award.
  \newblock \url{https://rohan.padhye.org/files/datetimebugs-msr25.pdf}

\bibitem{liva2018}
  Giovanni Liva.
  \newblock \emph{Modeling Time in Java Programs for Automatic Error
  Detection}.
  \newblock PhD thesis, Alpen-Adria-Universit\"at Klagenfurt, 2018.

\bibitem{bry-spranger-multicalendar}
  Fran\c{c}ois Bry and Stephanie Spranger.
  \newblock Towards a multi-calendar temporal type system for
  (Semantic) Web query languages.
  \newblock In \emph{Principles and Practice of Semantic Web
  Reasoning (PPSWR 2004)}, volume 3208 of \emph{Lecture Notes in
  Computer Science}, pages 90--104. Springer, 2004.
  \newblock \doi{10.1007/978-3-540-30122-6\_8}

\bibitem{gruninger2017allenontology}
  Michael Gr\"uninger and Zhuojun Li.
  \newblock The time ontology of Allen's interval algebra.
  \newblock In \emph{24th International Symposium on Temporal
  Representation and Reasoning (TIME 2017)}, volume 90 of
  \emph{LIPIcs}, pages 16:1--16:16. Schloss
  Dagstuhl -- Leibniz-Zentrum f\"ur Informatik, 2017.
  \newblock \doi{10.4230/LIPIcs.TIME.2017.16}

\bibitem{dyreson2021logsegmented}
  Curtis E. Dyreson and M. A. Manazir Ahsan.
  \newblock Achieving a sequenced, relational query language with
  log-segmented timestamps.
  \newblock In \emph{28th International Symposium on Temporal
  Representation and Reasoning (TIME 2021)}, volume 206 of
  \emph{LIPIcs}, pages 14:1--14:13. Schloss
  Dagstuhl -- Leibniz-Zentrum f\"ur Informatik, 2021.
  \newblock \doi{10.4230/LIPIcs.TIME.2021.14}

\bibitem{dyreson2023logsegmented}
  Curtis E. Dyreson.
  \newblock Optimization of nonsequenced queries using log-segmented
  timestamps.
  \newblock In \emph{30th International Symposium on Temporal
  Representation and Reasoning (TIME 2023)}, volume 278 of
  \emph{LIPIcs}, pages 13:1--13:15. Schloss
  Dagstuhl -- Leibniz-Zentrum f\"ur Informatik, 2023.
  \newblock \doi{10.4230/LIPIcs.TIME.2023.13}

\end{thebibliography}



\end{document}
